getEnlargement() --> O(1) as computes one enlarged rectangle and subtract two areas.

computeMBR() --> loops over all items or children in a single node. A node holds at most maxChildren entries.
        O(maxChildren)  →  treated as O(1) since maxChildren is a fixed constant

splitNode() --> PickSeeds step compares every pair of entries in the overflowed node, which has maxChildren + 1 entries. The PickNext step assigns each remaining entry in one pass.
                so PickSeeds:O(maxChildren)
                PickNext: O(maxChildren)
                Total O(2 x maxChildren) →  O(1) since maxChildren is fixed

adjustTree() --> walks from leaf to root , one level at a time. Haight is log n , work at each level is O(1).
                    so, O(logn)

insert() --> calls chooseleaf+ possibly splitNode + adjustTree
            so, chooseLeaf:   O(log n)
            splitNode:    O(1)  at most one split per level
            adjustTree:   O(log n)
            Total:        O(log n)

findleaf() --> in the worst case it must visit every node in the tree because MBRs overlap and multiple branches must be explored. 
                In practice with low overlap it behaves like O(log n), but worst case it touches every node.

search() --> Low overlap means few branches. High overlap means many branches.
            Best case:   O(log n + k)   where k = number of results returned
            Worst case:  O(n)

condenseTree() -->  Walk up:      O(log n)
                    Re-insertions: O(m x log n)
                    Total:        O(m x log n) ; If d nodes are dissolved with m total orphaned entries

remove() --> findLeaf:      O(log n) average, O(n) worst
            condenseTree:  O(log n) average
            Total:         O(log n) average,  O(n) worst case